% !TeX program = pdflatex
\documentclass[11pt,a4paper]{amsart}

\newcommand{\notelanguage}{finnish}
% Shared preamble for course notes and exercise sets.

\usepackage[T1]{fontenc}
\usepackage{lmodern}
\usepackage[english]{babel}

\usepackage[a4paper,margin=30mm]{geometry}
\usepackage{microtype}
\usepackage{graphicx}
\usepackage{enumitem}

\usepackage{amsmath}
\usepackage{amssymb}
\usepackage{amsthm}
\usepackage{mathtools}
\usepackage{aliascnt}

\usepackage[hidelinks,pdfusetitle]{hyperref}

\numberwithin{equation}{section}
\setlist{itemsep=0.25em,topsep=0.5em}

\theoremstyle{plain}
\newtheorem{theorem}{Theorem}[section]
\newaliascnt{lemma}{theorem}
\newtheorem{lemma}[lemma]{Lemma}
\aliascntresetthe{lemma}
\newaliascnt{proposition}{theorem}
\newtheorem{proposition}[proposition]{Proposition}
\aliascntresetthe{proposition}
\newaliascnt{corollary}{theorem}
\newtheorem{corollary}[corollary]{Corollary}
\aliascntresetthe{corollary}

\theoremstyle{definition}
\newaliascnt{definition}{theorem}
\newtheorem{definition}[definition]{Definition}
\aliascntresetthe{definition}
\newaliascnt{example}{theorem}
\newtheorem{example}[example]{Example}
\aliascntresetthe{example}
\newaliascnt{problem}{theorem}
\newtheorem{problem}[problem]{Problem}
\aliascntresetthe{problem}

\theoremstyle{remark}
\newaliascnt{remark}{theorem}
\newtheorem{remark}[remark]{Remark}
\aliascntresetthe{remark}

\usepackage[nameinlink,noabbrev]{cleveref}

\crefname{theorem}{theorem}{theorems}
\crefname{lemma}{lemma}{lemmas}
\crefname{proposition}{proposition}{propositions}
\crefname{corollary}{corollary}{corollaries}
\crefname{definition}{definition}{definitions}
\crefname{example}{example}{examples}
\crefname{problem}{problem}{problems}
\crefname{remark}{remark}{remarks}

\DeclareMathOperator{\im}{im}
\DeclareMathOperator{\rank}{rank}
\DeclarePairedDelimiter{\abs}{\lvert}{\rvert}
\DeclarePairedDelimiter{\norm}{\lVert}{\rVert}
\DeclarePairedDelimiter{\set}{\lbrace}{\rbrace}

\usepackage{tikz}

\title{Analyysi I\\Harjoituksen 3 ratkaisut}
\author{}
\date{}

\begin{document}

\exerciseheading{Analyysi I}{Harjoituksen 3 ratkaisut}

\begin{exercise}
  Olkoon \((x_n)\) lukujono, jolle \(x_n = \frac{3n}{n + 37}\). Osoita suoraan
  määritelmästä, että lukujono suppenee kohti lukua \(3\). Osoita lisäksi,
  että lukujono on kasvava.
\end{exercise}

\begin{solution}
\end{solution}

\begin{exercise}
  \leavevmode\par
\begin{enumerate}[label=\alph*), nosep]
  \item Olkoon \((x_n)\) suppeneva lukujono. Määritellään lukujono \((y_n)\)
    asettamalla, että \(y_n = \frac{1}{\sqrt{n}}x_n\). Osoita, että lukujono
    \((y_n)\) suppenee.
  \item Oletetaan, että lukujono \((x_n)\) on rajoitettu. Määritellään lukujono
    \((y_n)\) kuten (a)-kohdassa. Mitä voidaan sanoa lukujonon \((y_n)\)
    suppenemisesta?
\end{enumerate}
\end{exercise}

\begin{solution}
\end{solution}

\begin{exercise}
Anna esimerkki lukujonosta, joka
\begin{enumerate}[label=\alph*), nosep]
\item on kasvava, mutta ei aidosti kasvava, ja hajaantuu.
\item on aidosti vähenevä ja suppenee.
\end{enumerate}
\end{exercise}

\begin{solution}
\end{solution}

\begin{exercise}
Olkoon \(E\) reaalilukujen \(\RR\) epätyhjä alhaalta rajoitettu osajoukko.
Osoita, että jokaista \(\epsilon > 0\) kohti on olemassa sellainen \(x \in E\),
että \(\inf E \leq x < \inf E + \epsilon\).
\end{exercise}

\begin{solution}
\end{solution}

\begin{exercise}
Tutustu kirjan Lauseeseen 2.3.4 (Bolzanon-Weierstrassin lause) ja sen
todistukseen. Bolzanon-Weierstrassin lauseen todistuksessa käytettiin joukkojen
\(\set{x_n \colon n \geq k}\) supremumeja. Muuta todistusta niin että se käyttää
samojen joukkojen infimumeja.
\end{exercise}

\begin{solution}
\end{solution}

\begin{exercise}
Sinifunktion keskeinen ominaisuus on, että sen arvot kuuluvat välille
\([-1, 1]\) eli \(\sin(x) \in [-1, 1]\) kaikilla \(x \in \RR\). Määritellään
lukujono \((x_n)\) asettamalla, että
\[
x_n = \frac{12\sin(\frac{3\pi n}{8})^4}{n(2 + \sin(\frac{\pi n}{2}))}
\]
kaikilla \(n \in \NN\). Osoita suoraan määritelmästä, että lukujono \((x_n)\)
suppenee kohti lukua \(0\).
\end{exercise}

\begin{solution}
\end{solution}
\end{document}
